\subsection{Uniform Limits}

% **Role:** convergence of functions and analytic approximation.
%
% Pointwise versus uniform convergence, uniform metric, M-test, preservation of continuity, differentiation, integration.


Uniform convergence asks when limits of functions preserve structure.\\
Whether $N$ may depend on $x$ controls continuity, integration, differentiation, and the supremum metric.

\begin{perspective} \label{per:function-limits-before-function-sequences} \hlt{Function Limits Before Function-Sequence Limits}\\
$\lim\limits_{t\to x}f(t)$ denotes the metric-space function limit of Definition \ref{def:function-limit-metric}.\\
For complex-valued functions with individual function limits, uniform convergence permits the limits in $n$ and $t$ to commute.
\end{perspective}

\begin{definition} \label{def:pointwise-convergence-functions} \hlt{Pointwise Limits of Functions}\\
Let $(Y,d_Y)$ be a metric space and let $f_n,f:E\to Y$.
\begin{enumerate}[label=(\roman*)]
\setlength{\itemsep}{0pt}
\item $f_n\to f$ \hlt{pointwise} on $E$ if $(\forall x\in E)\ d_Y(f_n(x),f(x))\to0$.
\item If $Y\in\{\R,\C\}$, set $s_N(x)\defeq\sum\limits_{n=1}^{N}f_n(x)$. Then $\sum f_n\to f$ \hlt{pointwise} on $E$ if $(\forall x\in E)\ s_N(x)\to f(x)$.
\end{enumerate}
No metric on $E$ is required; topology matters only for preserving continuity, limits, or integrals.
\end{definition}

\begin{example} \label{ex:pointwise-limits-do-not-preserve-structure} \hlt{Pointwise Limits Can Lose Structure}
\begin{enumerate}[label=(\roman*)]
\setlength{\itemsep}{0pt}
\item For $m,n\in\N$, $s_{m,n}\defeq\frac{m}{m+n}$ satisfies $\lim\limits_{n\to\infty}\lim\limits_{m\to\infty}s_{m,n}=1$ but $\lim\limits_{m\to\infty}\lim\limits_{n\to\infty}s_{m,n}=0$.
\item On $\R$, set $f_n(x)=\frac{x^2}{(1+x^2)^n}$. Then $\sum\limits_{n=0}^{\infty}f_n$ equals $0$ at $0$ and $1+x^2$ otherwise, by Theorem \ref{thm:geometric-series-nonnegative}.
\item On $\R$, $f_n(x)=\frac{\sin(nx)}{\sqrt n}$ satisfies $f_n\to0$ pointwise but $f_n'(0)=\sqrt n\to+\infty$.
\item $f_n(x)=n^2x(1-x^2)^n$ on $[0,1]$ satisfies $f_n\to0$ pointwise, while $\int_0^1 f_n\,dx=\frac{n^2}{2n+2}\to+\infty$.
\end{enumerate}
\hlt{Key Insight}: Pointwise control can miss global spikes and oscillations.
\end{example}

\begin{definition} \label{def:uniform-convergence-functions} \hlt{Uniform Convergence}\\
Let $(Y,d_Y)$ be a metric space and let $f_n,f:E\to Y$.\\
Then $f_n\rightrightarrows f$ on $E$ if $(\forall\varepsilon>0)(\exists N\in\N)(\forall n\geq N)(\forall x\in E)\ d_Y(f_n(x),f(x))\leq\varepsilon$.\\
In the scalar-valued case, take $d_Y(u,v)=\abs{u-v}$; $\sum f_n$ converges uniformly on $E$ if its partial sums do.
\end{definition}

\begin{perspective} \label{per:uniform-convergence-epsilon-strip} \hlt{One Tail, All Points}\\
Pointwise convergence lets $N$ depend on $x$; uniform convergence makes $N$ independent of $x$.\\
Equivalently, late graphs lie in one $\varepsilon$-strip around $f$, as illustrated in Figure \ref{fig:uniform-convergence-epsilon-strip}.
\end{perspective}

\begin{illustrationfigure}[H]
  \centering
  \begin{tikzpicture}[
  >=Stealth,
  every node/.style={font=\small},
  axis/.style={->, gray!70, line width=0.45pt},
  strip/.style={gray!65, dotted, line width=0.65pt},
  limit/.style={black, line width=0.9pt},
  tail/.style={blue!70!black, line width=0.75pt},
  domainbar/.style={black!75, line width=1.3pt}
]
  \draw[axis] (-0.35,0) -- (7.2,0) node[right] {$x$};
  \draw[axis] (0,-0.5) -- (0,3.3) node[above] {$y$};

  \draw[domainbar] (0.75,-0.18) -- (6.45,-0.18);
  \node[below] at (3.6,-0.18) {$E$};

  \draw[strip, domain=0.75:6.45, variable=\xx, samples=120]
    plot (\xx,{2.0 + 0.24*sin(85*\xx) + 0.52});
  \draw[limit, domain=0.75:6.45, variable=\xx, samples=120]
    plot (\xx,{2.0 + 0.24*sin(85*\xx)});
  \draw[strip, domain=0.75:6.45, variable=\xx, samples=120]
    plot (\xx,{2.0 + 0.24*sin(85*\xx) - 0.52});

  \draw[tail, domain=0.75:6.45, variable=\xx, samples=120]
    plot (\xx,{2.0 + 0.24*sin(85*\xx) + 0.18*sin(190*\xx)});
  \draw[tail, blue!45!black, domain=0.75:6.45, variable=\xx, samples=120]
    plot (\xx,{2.0 + 0.24*sin(85*\xx) - 0.15*cos(170*\xx)});

  \node[right] at (6.52,2.78) {$f+\varepsilon$};
  \node[right] at (6.52,2.18) {$f$};
  \node[right] at (6.52,1.68) {$f-\varepsilon$};
  \node[blue!65!black, anchor=west] at (1.25,2.95) {$f_n,\ n\geq N$};
\end{tikzpicture}

  \caption{One index $N$ puts all later graphs inside the same $\varepsilon$-strip around $f$ on $E$.}
  \label{fig:uniform-convergence-epsilon-strip}
\end{illustrationfigure}

\begin{example} \label{ex:xn-pointwise-not-uniform} \hlt{$x^n$ Converges Pointwise But Not Uniformly}\\
On $[0,1]$, let $g_n(x)=x^n$. Then $g_n\to g$ pointwise, where $g(x)=
\begin{cases}
0, & 0\leq x<1,\\
1, & x=1.
\end{cases}$.\\
The convergence is not uniform: $(\forall n\geq1)(\exists x\in[0,1))\ \abs{x^n-g(x)}=x^n>\frac{1}{2}$.\\
Figure \ref{fig:xn-pointwise-not-uniform} plots the dimensionless functions $x,x^2,x^4,x^8$ over $x\in[0,1]$ and exposes the persistent endpoint transition.
\end{example}

\begin{visualizationfigure}[H]
  \centering
  \begin{tikzpicture}[
  >=Stealth,
  every node/.style={font=\small},
  axis/.style={->, gray!70, line width=0.45pt},
  curve/.style={line width=0.75pt},
  limit/.style={gray!70, dashed, line width=0.55pt},
  point/.style={circle, fill=black, inner sep=1.35pt}
]
  \draw[axis] (-0.15,0) -- (4.7,0) node[right] {$x$};
  \draw[axis] (0,-0.15) -- (0,3.0) node[above] {$y$};

  \draw[gray!65] (4,0.05) -- (4,-0.05) node[below] {$1$};
  \draw[gray!65] (0.05,2.4) -- (-0.05,2.4) node[left] {$1$};

  \draw[curve, blue!70!black, domain=0:4, variable=\xx, samples=80] plot (\xx,{2.4*(\xx/4)});
  \draw[curve, dartmouthgreen!75!black, domain=0:4, variable=\xx, samples=100] plot (\xx,{2.4*(\xx/4)^2});
  \draw[curve, jazzberryjam!80!black, domain=0:4, variable=\xx, samples=100] plot (\xx,{2.4*(\xx/4)^4});
  \draw[curve, black!70, domain=0:4, variable=\xx, samples=120] plot (\xx,{2.4*(\xx/4)^8});

  \draw[limit] (0,0) -- (4,0);
  \node[point] at (4,2.4) {};
  \node[right] at (4.1,2.4) {$g(1)=1$};
  \node[below, black!70] at (2.15,-0.02) {$g(x)=0$ for $x<1$};

  \node[blue!70!black] at (2.1,1.55) {$x$};
  \node[dartmouthgreen!75!black] at (2.55,0.82) {$x^2$};
  \node[jazzberryjam!80!black] at (3.08,0.55) {$x^4$};
  \node[black!70] at (3.5,0.32) {$x^8$};
\end{tikzpicture}

  \caption{$x^n\to0$ pointwise on $[0,1)$, but the endpoint prevents uniform convergence on $[0,1]$.}
  \label{fig:xn-pointwise-not-uniform}
\end{visualizationfigure}

\begin{theorem} \label{thm:uniform-cauchy-criterion} \hlt{Cauchy Criterion for Uniform Convergence}\\
For $f_n:E\to\C$, uniform convergence on $E$ is equivalent to $(\forall\varepsilon>0)(\exists N\in\N)(\forall m,n\geq N)(\forall x\in E)\ \abs{f_n(x)-f_m(x)}\leq\varepsilon$.
\begin{proof}
\begin{enumerate}[label=\roman*.]
\setlength{\itemsep}{0pt}
\item[$(\Rightarrow)$] Fix $\varepsilon>0$. If $f_n\rightrightarrows f$, choose $N$ so that $(k\geq N,\ x\in E)\Rightarrow\abs{f_k(x)-f(x)}\leq\frac{\varepsilon}{2}$. Then $(m,n\geq N,\ x\in E)\Rightarrow\abs{f_n(x)-f_m(x)}\leq\varepsilon$.
\item[$(\Leftarrow)$] For $x\in E$, $(f_n(x))$ is Cauchy in $\C\cong\R^2$ by the stated condition, so it converges by Corollary \ref{cor:euclidean-cauchy-convergence}.\\
Set $f(x)\defeq\lim\limits_{j\to\infty}f_j(x)$. Fix $\varepsilon>0$ and choose $N$ so that $(m,n\geq N,\ x\in E)\Rightarrow\abs{f_n(x)-f_m(x)}\leq\frac{\varepsilon}{2}$.\\
For $n\geq N$ and $x\in E$, choose $m\geq N$ with $\abs{f_m(x)-f(x)}<\frac{\varepsilon}{2}$ by $f_j(x)\to f(x)$. Then $\abs{f_n(x)-f(x)}\leq\varepsilon$, so $f_n\rightrightarrows f$.
\end{enumerate}
\hlt{Structural Commentary}: The same $N$ controls the pointwise Cauchy tails simultaneously across $E$.\\
\hlt{Key Insight}: One Cauchy index controls every point of $E$.
\end{proof}
\end{theorem}

\begin{theorem} \label{thm:uniform-convergence-supremum-test} \hlt{Supremum Test for Uniform Convergence}\\
Let $\emptyset\ne E$, $f_n,f:E\to\C$. Set $M_n\defeq\sup\limits_{\overline{\R}}\{\abs{f_n(x)-f(x)}:x\in E\}\in[0,+\infty]$. Then $f_n\rightrightarrows f$ on $E$ if and only if $M_n\to0$.
\begin{proof}
\begin{enumerate}[label=\roman*.]
\setlength{\itemsep}{0pt}
\item[$(\Rightarrow)$] Fix $\varepsilon>0$. Choose $N$ with $(\forall n\geq N)(\forall x\in E)\ \abs{f_n(x)-f(x)}\leq\varepsilon$ by $f_n\rightrightarrows f$. Then $(\forall n\geq N)\ M_n\leq\varepsilon$.
\item[$(\Leftarrow)$] Fix $\varepsilon>0$. Choose $N$ with $(\forall n\geq N)\ M_n\leq\varepsilon$ by $M_n\to0$. Then $(\forall n\geq N)(\forall x\in E)\ \abs{f_n(x)-f(x)}\leq\varepsilon$.
\end{enumerate}
\hlt{Structural Commentary}: Extended suprema package every pointwise error into the scalar sequence $(M_n)$.\\
\hlt{Key Insight}: The global error vanishes exactly when one index works everywhere.
\end{proof}
\end{theorem}

\begin{theorem} \label{thm:weierstrass-m-test} \hlt{Weierstrass $M$-Test}\\
Let $f_n:E\to\C$, $(M_n)\subseteq[0,+\infty)$ satisfy $(\forall n\in\N)(\forall x\in E)\ \abs{f_n(x)}\leq M_n$. If $\sum M_n$ converges, then $\sum f_n$ converges uniformly on $E$.
\begin{proof}
Set $s_j\defeq\sum\limits_{k=1}^{j}f_k$. If $m>n$ and $x\in E$, then $\abs{s_m(x)-s_n(x)}\leq\sum\limits_{k=n+1}^{m}M_k$. Fix $\varepsilon>0$.\\
Choose $N\in\N$ s.t. $(\forall m>n\geq N)\ \sum\limits_{k=n+1}^{m}M_k<\varepsilon$ by Theorem \ref{thm:series-cauchy-criterion}. For $m,n\geq N$, interchange indices if necessary so that $m\geq n$.\\
Difference is $0$ if $m=n$, preceding bound less than $\varepsilon$ if $m>n$; $(s_j)$ is uniformly Cauchy, converges uniformly by Theorem \ref{thm:uniform-cauchy-criterion}.\\
\hlt{Structural Commentary}: Triangle inequality transfers numerical Cauchy tails to uniform bounds for the partial sums.\\
\hlt{Key Insight}: Summable scalar majorants dominate every functional tail.
\end{proof}
\end{theorem}

\begin{theorem} \label{thm:uniform-convergence-interchanges-function-limits} \hlt{Uniform Convergence Interchanges Limits}\\
Let $(X,d)$ be metric space, $E\subseteq X$, $p$ be limit point of $E$, $f_n,f:E\to\C$ satisfy $f_n\rightrightarrows f$ on $E$. Suppose $A_n\defeq \lim\limits_{t\to p}f_n(t)\in\C$ for $n\in\N$.
\begin{enumerate}[label=(\roman*)]
\setlength{\itemsep}{0pt}
\item $(A_n)$ converges;
\item $\lim\limits_{t\to p} f(t)=\lim\limits_{n\to\infty}A_n$.
\end{enumerate}
\begin{proof}
\begin{enumerate}[label=(\roman*)]
\setlength{\itemsep}{0pt}
\item Fix $\varepsilon>0$ and choose $N$ with $(\forall m,n\geq N)(\forall t\in E)\ \abs{f_n(t)-f_m(t)}\leq\frac{\varepsilon}{3}$ by Theorem \ref{thm:uniform-cauchy-criterion}.\\
For $m,n\geq N$, $\delta>0$, $(t\in E,\ 0<d(t,p)<\delta)\Rightarrow[\abs{f_n(t)-A_n}<\frac{\varepsilon}{3}\land\abs{f_m(t)-A_m}<\frac{\varepsilon}{3}]$ by definitions of $A_n$ and $A_m$.\\
$E\cap B_\delta(p)\setminus\{p\}\ne\emptyset$ by $p$ being a limit point of $E$; choose $t$ in this set; $\abs{A_n-A_m}<\varepsilon$ by triangle inequality.\\
Thus $(A_n)$ is Cauchy in $\C\cong\R^2$ and converges by Corollary \ref{cor:euclidean-cauchy-convergence}; set $A\defeq\lim\limits_{n\to\infty}A_n$.
\item For $n\in\N$, $\abs{f(t)-A}\leq \abs{f(t)-f_n(t)}+\abs{f_n(t)-A_n}+\abs{A_n-A}$.\\
Given $\varepsilon>0$, choose $n$ so that $(\forall t\in E)\ \abs{f(t)-f_n(t)}\leq\frac{\varepsilon}{4}$ and $\abs{A_n-A}<\frac{\varepsilon}{4}$.\\
Since $A_n=\lim\limits_{t\to p}f_n(t)$, choose $\delta>0$ so that $(t\in E,\ 0<d(t,p)<\delta)\Rightarrow\abs{f_n(t)-A_n}<\frac{\varepsilon}{4}$.\\
The three-term estimate is then less than $\frac{3\varepsilon}{4}$, so $\lim\limits_{t\to p}f(t)=A$.
\end{enumerate}
\hlt{Structural Commentary}: Uniform Cauchy control establishes $(A_n)\to A$; a three-term estimate identifies $\lim_{t\to p}f(t)=A$.\\
\hlt{Key Insight}: The shared error bound survives passage through the outer limit.
\end{proof}
\end{theorem}

\begin{theorem} \label{thm:uniform-convergence-preserves-continuity} \hlt{Uniform Limits Preserve Continuity}\\
Let $E$ be subspace of metric space, $f_n,f:E\to\C$. If $(f_n)$ consists of continuous functions, $f_n\rightrightarrows f$ on $E$, then $f$ continuous on $E$.
\begin{proof}
Let $x\in E$. If $x$ is isolated, $f$ is continuous at $x$. If $x$ is a limit point, then $(\forall n\in\N)\ \lim\limits_{t\to x}f_n(t)=f_n(x)$.\\
Then $\lim\limits_{t\to x}f(t)=\lim\limits_{n\to\infty}f_n(x)=f(x)$ by Theorem \ref{thm:uniform-convergence-interchanges-function-limits}.\\
\hlt{Structural Commentary}: The interchange theorem transfers each local limit identity from $f_n$ to $f$.\\
\hlt{Key Insight}: A common tail index preserves each local limit identity.
\end{proof}
\end{theorem}

\begin{theorem} \label{thm:dini-monotone-uniform-convergence} \hlt{Dini's Monotone Uniform Convergence Theorem}\\
Let $K$ be a compact metric space, $f_n,f:K\to\R$ be continuous. If $(\forall x\in K)\ f_n(x)\downarrow f(x)$, then $f_n\rightrightarrows f$ on $K$.
\begin{proof}
\begin{enumerate}[label=(\roman*)]
\setlength{\itemsep}{0pt}
\item \hlt{Bad sets}: Set $g_n\defeq f_n-f$. For $n\in\N$, $g_n$ is continuous by Theorem \ref{thm:continuous-algebra-and-components}(i).\\
The relation $(\forall x\in K)\ g_n(x)\downarrow0$ holds by hypothesis. Fix $\varepsilon>0$ and set $K_n\defeq g_n^{-1}([\varepsilon,+\infty))$.
\item \hlt{Compactness}: For $n\in\N$, $K_n$ is closed in $K$ by Corollary \ref{cor:continuity-closed-preimages}, hence compact by Theorem \ref{thm:closed-subsets-of-compact-sets-compact}.\\
Inclusions $K_{n+1}\subseteq K_n$ hold by $(\forall x\in K)\ g_{n+1}(x)\leq g_n(x)$. Identity $\bigcap\limits_{n=1}^{\infty}K_n=\emptyset$ holds by $(\forall x\in K)\ g_n(x)\to0$.\\
If $(\forall n\in\N)\ K_n\ne\emptyset$, then $\bigcap\limits_{n=1}^{\infty}K_n\ne\emptyset$ by Corollary \ref{cor:nested-nonempty-compact-intersection}, contradiction. Hence $(\exists N\in\N)\ K_N=\emptyset$.
\item \hlt{Uniform bound}: By nestedness, $(\forall n\geq N)(\forall x\in K)\ 0\leq f_n(x)-f(x)<\varepsilon$.
\end{enumerate}
\hlt{Structural Commentary}: Monotone convergence turns the bad sets into a nested compact sequence.\\
\hlt{Key Insight}: Compactness converts pointwise escape from nested bad sets into uniform escape.
\end{proof}
\end{theorem}

\begin{definition} \label{def:supremum-norm-continuous-functions} \hlt{The Supremum Metric on Bounded Continuous Functions}\\
Let $X$ be a nonempty metric space.
\begin{enumerate}[label=(\roman*)]
\setlength{\itemsep}{0pt}
\item $C_b(X)$ is the set of bounded continuous maps $X\to\C$.
\item For $f\in C_b(X)$, $\norm{f}_{\infty}\defeq\sup\limits_{x\in X}\abs{f(x)}$.
\item The \hlt{supremum metric} on $C_b(X)$ is $d_{\infty}(f,g)\defeq\norm{f-g}_{\infty}$.
\item For $(f_n)\subseteq C_b(X)$ and $f\in C_b(X)$, $f_n\to f$ in $d_{\infty}\Longleftrightarrow\norm{f_n-f}_{\infty}\to0\Longleftrightarrow f_n\rightrightarrows f$ on $X$.
\end{enumerate}
\hlt{Key Insight}: Uniform convergence is metric convergence in $C_b(X)$.
\end{definition}

\begin{perspective} \label{per:supremum-metric-function-space-lens} \hlt{A Metric Space Whose Points Are Functions}\\
For a nonempty metric space $X$, uniform convergence in $C_b(X)$ is sequence convergence in $(C_b(X),d_{\infty})$.
\end{perspective}

\begin{perspective} \label{per:four-limit-notions-function-space-view} \hlt{Four Limit Notions}\\
The four limit notions are sequence limits, function limits, pointwise function limits, and uniform function limits.\\
For a nonempty metric space $X$, $(f_n)\subseteq C_b(X)$, and $f\in C_b(X)$, the last is $d_{\infty}$-convergence by Definition \ref{def:supremum-norm-continuous-functions}(iv).
\end{perspective}
